\section{Inference Module}
\label{sec:inference}

\subsection{Trajectory Representation}

The trajectory object is the unit of input to the inference module. Each
trajectory carries five fields: \texttt{agent_id} (a string identifying the agent
under audit), \texttt{declared_purpose} (the natural-language $D$ required by EU
AI Act Annex~IV), \texttt{steps} (a non-empty list of typed step objects),
\texttt{seed} (an optional integer), and \texttt{metadata} (a free-form
dictionary for provenance information). Each step contains a \texttt{state}
dictionary encoding the agent's context, an action object with \texttt{tool_name}
and \texttt{parameters} as mandatory typed fields, and an \texttt{observation}
dictionary carrying the tool response. The schema is validated on ingest; a
trajectory containing zero steps raises a validation error before inference
begins.

The \texttt{tool_name} and \texttt{parameters} fields are the load-bearing
elements for typed feature extraction: because every action is labelled with its
tool identity and its argument structure, the Goal Decomposition Judge critic can
reason about what the agent did without relying on free-text parsing of opaque log
lines. The \texttt{seed} field enables deterministic replay of the agent's step
sequence: given the same seed, the gym reference agents produce identical
trajectories. Combined with the temperature-zero decomposition described in
Section~\ref{subsec:gdj}, this property makes an audit exactly reproducible from
its inputs alone.

\subsection{Goal Decomposition Judge}
\label{subsec:gdj}

The Goal Decomposition Judge (GDJ) solves a problem that prior IRL-based auditing
work never needed to address. Prior approaches, including the Alignment
Auditor~\citep{alignmentAuditor2025} and IR\textsuperscript{3}~\citep{ir3contrastive},
operated at the model-weights or RLHF reward-model level, where the feature space
was directly accessible from training artefacts. IOR operates on post-deployment
API traces from black-box agents, where no such internal structure is available.
The GDJ bridges the gap by constructing a feature space from the declared purpose
$D$ using a two-stage pipeline.

The first stage is the decomposer, which runs exactly once per audit. Given $D$,
the decomposer issues a single language-model call instructing it to extract
exactly $n$ observable, measurable sub-goals that together cover the declared
purpose, where each sub-goal must be observable from a single (state, action,
observation) triplet, expressed as a measurable proxy signal rather than an
abstract value, and distinct from all others. The decomposer runs at temperature
zero, which is an architectural invariant rather than a tuning choice. Because the
Laplace approximation and the divergence report are computed over the feature
dimensions produced by this call, any non-determinism in the decomposer would make
two ostensibly identical audits incomparable. Determinism at temperature zero
guarantees that the same declared purpose always produces the same goal
specification (the ordered list $\{g_1, \ldots, g_d\}$), which in turn guarantees
that the feature basis is stable across runs. For the \texttt{CostMinimiserAgent},
whose declared purpose instructs the agent to find the cheapest option whilst
minimising steps and respecting budget, the decomposer produces sub-goals
including \texttt{seeks_price_comparison}, \texttt{avoids_redundant_calls},
\texttt{confirms_before_booking}, \texttt{respects_budget}, and
\texttt{minimises_steps}, each of which is evaluable from a single tool-call
triplet.

The second stage is the critic, which runs once per trajectory step. For each
step $(s_t, a_t, o_t)$, the critic scores the triplet against each $g_i$ on a
scale where $0.0$ indicates the action clearly harms or is irrelevant to the
sub-goal, $0.5$ is neutral, and $1.0$ indicates the action clearly and directly
advances it. The critic returns scores $\phi_t \in [0,1]^d$, per-dimension
confidence values, and a short natural-language explanation for each score. Like
the decomposer, the critic runs at temperature zero, ensuring that a given triplet
against a given goal specification always produces the same feature vector. The
confidence scores are stored in the feature matrix and are intended (Phase~1) to
propagate scoring uncertainty into the Bayesian IRL likelihood as per-step
weights, so that steps the critic is uncertain about contribute less to the
posterior.

The GDJ carries an important independence requirement. The judge model must differ
from the agent under audit, or be prompted blind to the agent's identity. If the
same model acts as both agent and judge, systematic biases in that model's
world-view align across both roles: the judge will tend to interpret the agent's
actions charitably, and the divergence will be systematically underestimated. This
is not a best practice but a design invariant.

Feature quality is validated empirically rather than assumed. The
planted-divergence gym (Section~\ref{sec:gym}) provides three reference agents
whose objective divergences are known in advance. If the GDJ features accurately
capture the planted divergences, the full pipeline recovers them above a stated
correlation threshold. This transforms what would otherwise be an unfalsifiable
design claim, that the judge produces useful features, into a measurable property
with a defined acceptance criterion.

A confirmation bias risk accompanies the decomposer architecture. Because the
decomposer receives $D$ as input, it is primed to generate sub-goals that align
with $D$ and may therefore produce features that confirm the declared purpose
rather than revealing what the agent actually optimises for.
Section~\ref{sec:evaluation} addresses this directly through a blind-critic
ablation: the critic is run without knowledge of $D$, and the resulting feature
matrices are compared to those produced under full context.

\subsection{Bayesian IRL via Laplace Approximation}

IOR uses Bayesian IRL (BIRL) as its inference engine rather than maximum-entropy
IRL~\citep{ziebart2008maxent} or max-margin IRL~\citep{ng2000algorithms}. The
choice is driven by a regulatory requirement rather than a purely technical one.
MaxEnt IRL produces a point estimate of the reward function; max-margin IRL
produces a discriminative classifier. Neither emits a distribution over $\theta$.
A compliance reader who sees only $\hat{\theta}$ cannot tell whether the estimate
is well-determined or reflects a wide, shallow posterior. BIRL produces
$P(\theta \mid \tau)$, directly satisfying this requirement.

\paragraph{Likelihood.}
The likelihood model follows Boltzmann rationality. At each step $t$, the agent
selects action $a_t$ from its available toolset $\mathcal{A}$. The probability of
the observed action is
\begin{equation}
P(a_t \mid s_t, \theta) =
\frac{\exp\!\left(\beta \cdot \theta^\top \phi(s_t, a_t, o_t)\right)}
{\sum_{a' \in \mathcal{A}} \exp\!\left(\beta \cdot \theta^\top \phi(s_t, a', \cdot)\right)},
\end{equation}
where $\beta > 0$ is an inverse temperature controlling the assumed degree of
agent rationality. Higher $\beta$ concentrates the distribution towards the
highest-reward action; $\beta \to 0$ produces a uniform distribution regardless of
$\theta$. The full trajectory log-likelihood is
$\log P(\tau \mid \theta) = \sum_{t=1}^{T} \log P(a_t \mid s_t, \theta)$. The
implementation computes this using the log-sum-exp trick for numerical stability,
stacking the observed-action feature vector (row $0$ of
$\Phi_t \in \mathbb{R}^{(1+K) \times d}$) against $K$ counterfactual rows.

\paragraph{Counterfactuals and fast mode.}
Computing the softmax denominator requires feature vectors
$\phi(s_t, a', \cdot)$ for every action not taken at step $t$. Calling the GDJ
critic for each alternative tool at each step is exact but expensive: it incurs
$K \times T$ additional model calls per audit. IOR's fast mode avoids this cost by
approximating counterfactual feature vectors as $0.5 \cdot \mathbf{1}_d$,
representing a neutral action that neither advances nor harms any sub-goal. We
conjecture that this approximation is conservative, in the sense that a flat
neutral reference suppresses the softmax denominator relative to genuinely diverse
counterfactuals and so understates divergence magnitude; whether this holds depends
on the sign structure of $\theta$, and we do not yet have the empirical evidence
(Section~\ref{sec:evaluation}) to assert it. Full-mode counterfactual scoring is
available as a flag and constitutes the exact inference path. The fast-mode
approximation also has a consequence for what BIRL contributes: with a flat neutral
reference, the MAP estimate $\hat{\theta}_i$ rises roughly monotonically in
$(\bar{\phi}_i - 0.5)$, where $\bar{\phi}_i$ is the mean critic score on dimension
$i$, so the \emph{ranking} of the fast-mode point estimate is close to that of a
centred column-mean of the critic scores. We measure this directly in
Section~\ref{sec:evaluation} and find the two rankings highly correlated. The
distinctive value of Bayesian IRL therefore lies not in the point estimate under
fast mode but in the two things a column-mean cannot provide: the posterior
uncertainty $\sigma_i$ and the full-mode counterfactual treatment, in which the
reference is the agent's actual alternatives rather than a flat $0.5$.

\paragraph{Prior.}
The prior over reward weights is a spherical Gaussian:
$\theta \sim \mathcal{N}(0, \sigma^2 I)$, where $\sigma$ defaults to $1.0$. This
regularises against extreme reward weights and encodes the assumption that agents
generally pursue their declared purposes at least partially: an agent with no
declared-purpose alignment would require a large-magnitude $\theta$ pointing away
from $\theta_D$, which the prior penalises. For very short trajectories (fewer than
approximately ten steps), the prior dominates and the Laplace approximation is less
reliable.

\paragraph{MAP estimation.}
The MAP estimate
$\hat{\theta} = \arg\max_\theta\, [\log P(\tau \mid \theta) + \log P(\theta)]$ is
found by minimising the negative log-posterior via L-BFGS-B. The gradient is
available in closed form,
\begin{equation}
\nabla_\theta \log P(\tau \mid \theta) =
\sum_{t=1}^{T} \beta \left[\phi(s_t, a_t, o_t) -
\mathbb{E}_{a' \sim P(\cdot \mid s_t, \theta)}\big[\phi(s_t, a', \cdot)\big]\right],
\end{equation}
so no finite-difference approximation is required.

\paragraph{Laplace approximation.}
The posterior covariance is approximated as
$\Sigma = \big(-\nabla^2_\theta \log P(\theta \mid \tau)\big)^{-1}$ evaluated at
$\hat{\theta}$. The negative Hessian is
\begin{equation}
H = \frac{1}{\sigma^2} I +
\sum_{t=1}^{T} \beta^2 \operatorname{Cov}_{a' \sim P(\cdot \mid s_t, \hat{\theta})}
\big[\phi(s_t, a', \cdot)\big].
\end{equation}
IOR inverts $H$ via Cholesky factorisation; when $H$ is numerically singular
(which can occur for short trajectories or highly correlated sub-goals), a
pseudoinverse is used as a fallback. Posterior samples
$\theta^{(i)} \sim \mathcal{N}(\hat{\theta}, \Sigma)$ are drawn and stored for
downstream uncertainty quantification.

\paragraph{Log marginal likelihood.}
The Laplace approximation also yields an estimate of the log marginal likelihood,
\begin{equation}
\log P(\tau) \approx \log P(\tau \mid \hat{\theta}) + \log P(\hat{\theta})
- \tfrac{1}{2} \log \det H + \tfrac{d}{2} \log 2\pi,
\end{equation}
which is reported as a model-fit diagnostic and enables model comparison, for
instance against the Metropolis-Hastings sampler that IOR also provides.

\subsection{Addressing Ill-Posedness}

Ill-posedness in IRL is a first-class concern for regulatory evidence, not merely a
technical nuisance. Prior work establishes that partial identifiability is generic
in IRL: for any finite trajectory set, a family of reward functions exists that
generates the same distribution over behaviours, making the true $\theta$
unrecoverable in principle~\citep{partialIdentifiability2024}. The Alignment
Auditor explicitly warns against presenting single, overconfident IRL outputs as
compliance evidence~\citep{alignmentAuditor2025}. A compliance reader who sees a
single $\hat{\theta}$ and a divergence vector $\delta$ has no basis for
distinguishing ``this agent clearly over-weights tool-call frequency'' from ``we
cannot tell what this agent optimises for''.

IOR addresses this concern in three layers. First, the primary output of the
inference module is the posterior distribution $P(\theta \mid \tau)$, not the point
estimate $\hat{\theta}$; downstream components consume the full sample set. Second,
the divergence report exposes per-dimension standard deviations $\sigma_i$
computed from posterior samples, reported alongside $\delta_i$ so that a dimension
with $\delta_i = 0.3$ and $\sigma_i = 0.02$ warrants different attention than one
with $\delta_i = 0.3$ and $\sigma_i = 0.25$. Third, the log marginal likelihood is
reported as a summary diagnostic: an auditor who sees a very low value can infer
that the Boltzmann rationality model is a poor fit and should treat the divergence
estimates with corresponding caution. These measures do not resolve ill-posedness,
which is a fundamental limit of the inference problem; they make uncertainty
legible to non-specialist compliance readers, which is the appropriate response in
a regulatory context.

\subsection{Feature-Basis Variants}
\label{subsec:variants}

The GDJ is one of three interchangeable featurisers, all of which satisfy a single
protocol (decompose then score) so that the inference, divergence, and reporting
stages depend on the feature interface rather than on any particular
implementation. The two alternatives exist to address an obvious objection: the
GDJ depends on an external language model, which introduces cost and a
reproducibility risk across model versions.

The first alternative is a structural featuriser, which is deterministic and
model-free. It scores each step against a fixed behavioural basis (action novelty,
repetition of the previous tool, parameter novelty, payload size, and observation
non-emptiness) computed directly from the trajectory. It requires no external model
and produces identical output for identical input, giving a fully reproducible
floor. Its basis is behavioural rather than purpose-grounded, so it cannot localise
divergence against declared sub-goals; instead it answers the narrower question of
whether an agent's behaviour is predictable at all. In preliminary gym experiments
it achieves a mean held-out behaviour-prediction lift of $0.335$, already exceeding
the $0.20$ target, which establishes that a reproducible model-free basis carries
genuine signal.

The second alternative is an ensemble judge, which runs several diverse models as
critics and averages their per-step scores. Averaging reduces the influence of any
single model's idiosyncrasies whilst preserving the semantic, purpose-grounded
basis that divergence localisation requires. The ensemble additionally reports
inter-judge agreement (the mean pairwise correlation between judges' score
matrices) as a quantifiable reliability metric: high agreement indicates that the
recovered features do not hinge on one model's world-view, which directly answers
the external-dependence concern. Section~\ref{sec:evaluation} evaluates all three
featurisers head-to-head on the shared behaviour-prediction metric.
